% ====================================================================
% ФАЙЛ: tasks/03_electrodynamics/ohm_law_and_resistance_part1.tex
% ТЕМА: ЗАКОН ОМА ДЛЯ УЧАСТКА ЦЕПИ. УДЕЛЬНОЕ СОПРОТИВЛЕНИЕ
% ПАЧКА 1: ВАРИАНТЫ 1–10
% ====================================================================

% === ЗАДАЧА 1 ===
% Тег: #ЭЛЕКТРОДИНАМИКА #ТОК #ЗАКОН_ОМА #КИМ_11 #ВАРИАНТ_01
\begin{taskbasic}[code={\label{task:elem_ohm_v1_n11}}]
\textbf{Задача 1.} Во сколько раз увеличится сила тока, протекающего по металлический проводнику, если напряжение на его концах увеличить в $4$ раза, а длину проводника и площадь его поперечного сечения уменьшить в $2$ раза?

\kim{11}
\end{taskbasic}
\answer{4}

% === ЗАДАЧА 2 ===
% Тег: #ЭЛЕКТРОДИНАМИКА #ТОК #ЗАКОН_ОМА #КИМ_11 #ВАРИАНТ_02
\begin{taskbasic}[code={\label{task:elem_ohm_v2_n11}}]
\textbf{Задача 2.} Во сколько раз увеличится сила тока в цилиндрическом металлическом проводнике при постоянном напряжении на его концах, если его длину уменьшить в $2$ раза?

\kim{11}
\end{taskbasic}
\answer{2}

% === ЗАДАЧА 3 ===
% Тег: #ЭЛЕКТРОДИНАМИКА #ТОК #ЗАКОН_ОМА #КИМ_11 #ВАРИАНТ_05
\begin{taskbasic}[code={\label{task:elem_ohm_v5_n11}}]
\textbf{Задача 3.} На рисунке показаны графики зависимости силы тока $I$ от напряжения $U$ для двух резисторов. Чему равно отношение сопротивлений этих резисторов $\frac{R_1}{R_2}$?

\nopagebreak\smallskip
\begin{center}
\begin{tikzpicture}[scale=1.0, every node/.style={font=\footnotesize}]
    % Сетка gray!30
    \draw[xstep=1, ystep=1, gray!30, thin] (0,0) grid (4,4);
    
    % Оси координат
    \draw[-{Stealth[scale=0.8]}, black, thick] (0,0) -- (4.5,0) node[right] {$U$, В};
    \draw[-{Stealth[scale=0.8]}, black, thick] (0,0) -- (0,4.5) node[above] {$I$, А};
    
    % Графики VAH
    % Проводник 1: R1=2 Ом -> I = U/2 (проходит через (0,0), (2,1), (4,2))
    \draw[ultra thick, brandPrimary] (0,0) -- (4,2) node[right] {1};
    
    % Проводник 2: R2=1 Ом -> I = U/1 (проходит через (0,0), (2,2), (4,4))
    \draw[ultra thick, brandPrimary] (0,0) -- (4,4) node[above right] {2};
    
    % Точки на узлах
    \fill[black] (2,1) circle (1.5pt);
    \fill[black] (4,2) circle (1.5pt);
    \fill[black] (2,2) circle (1.5pt);
    \fill[black] (4,4) circle (1.5pt);
    
    % Пунктиры
    \draw[dashed, gray] (2,0) -- (2,2);
    \draw[dashed, gray] (4,0) -- (4,4);
    \draw[dashed, gray] (0,1) -- (2,1);
    \draw[dashed, gray] (0,2) -- (4,2);
    \draw[dashed, gray] (0,4) -- (4,4);
    
    % Оцифровка
    \node[left] at (0,1) {1};
    \node[left] at (0,2) {2};
    \node[left] at (0,4) {4};
    \node[below] at (2,0) {2};
    \node[below] at (4,0) {4};
    \node[below left] at (0,0) {0};
\end{tikzpicture}
\end{center}

\kim{11}
\end{taskbasic}
\answer{0,5}

% === ЗАДАЧА 4 ===
% Тег: #ЭЛЕКТРОДИНАМИКА #ТОК #ЗАКОН_ОМА #КИМ_11 #ВАРИАНТ_06
\begin{taskbasic}[code={\label{task:elem_ohm_v6_n11}}]
\textbf{Задача 4.} Медный провод длины $l = 100\text{ м}$ и площади поперечного сечения $S = 0{,}085\text{ мм}^2$ подключён к источнику постоянного напряжения. Чему равно электрическое сопротивление этого провода? (Удельное сопротивление меди $\rho = 0{,}017\text{ Ом}\cdot\text{мм}^2/\text{м}$). Ответ выразите в омах (Ом).

\kim{11}
\end{taskbasic}
\answer{20}

% ====================================================================
% ФАЙЛ: tasks/03_electrodynamics/ohm_law_and_resistance_part2.tex
% ТЕМА: ЗАКОН ОМА ДЛЯ УЧАСТКА ЦЕПИ. УДЕЛЬНОЕ СОПРОТИВЛЕНИЕ
% ПАЧКА 2: ВАРИАНТЫ 11–20
% ====================================================================

% === ЗАДАЧА 1 ===
% Тег: #ЭЛЕКТРОДИНАМИКА #ТОК #ЗАКОН_ОМА #КИМ_11 #ВАРИАНТ_11
\begin{taskbasic}[code={\label{task:elem_ohm_v11_n11}}]
\textbf{Задача 1.} Во сколько раз увеличится сила тока в цилиндрическом металлическом проводнике при постоянном напряжении на его концах, если площадь его поперечного сечения увеличить в $3$ раза?

\kim{11}
\end{taskbasic}
\answer{3}

% === ЗАДАЧА 2 ===
% Тег: #ЭЛЕКТРОДИНАМИКА #ТОК #ЗАКОН_ОМА #КИМ_11 #ВАРИАНТ_12
\begin{taskbasic}[code={\label{task:elem_ohm_v12_n11}}]
\textbf{Задача 2.} Имеется два цилиндрических проводника из одного и того же материала. Длина первого проводника в $3$ раза больше длины второго, а площадь поперечного сечения первого проводника в $2$ раза больше площади поперечного сечения второго. Чему равно отношение сопротивлений этих проводников $\frac{R_1}{R_2}$?

\kim{11}
\end{taskbasic}
\answer{1,5}

% === ЗАДАЧА 3 ===
% Тег: #ЭЛЕКТРОДИНАМИКА #ТОК #ЗАКОН_ОМА #КИМ_15 #ВАРИАНТ_17
\begin{taskbasic}[code={\label{task:elem_ohm_v17_n15}}]
\textbf{Задача 3.} На рисунке показаны графики зависимости силы тока $I$ от напряжения $U$ для двух цилиндрических проводников 1 и 2 одинаковой длины и одинаковой площади поперечного сечения, изготовленных из разных металлов. Из приведённого ниже списка выберите все верные утверждения относительно этих проводников.

\nopagebreak\smallskip
\begin{center}
\begin{tikzpicture}[scale=1.0, every node/.style={font=\footnotesize}]
    \draw[xstep=1, ystep=1, gray!30, thin] (0,0) grid (4,4);
    
    \draw[-{Stealth[scale=0.8]}, black, thick] (0,0) -- (4.5,0) node[right] {$U$, В};
    \draw[-{Stealth[scale=0.8]}, black, thick] (0,0) -- (0,4.5) node[above] {$I$, А};
    
    % График 1: R1 = 1 Ом (I = U)
    \draw[ultra thick, brandPrimary] (0,0) -- (4,4) node[above right] {1};
    
    % График 2: R2 = 2 Ом (I = U/2)
    \draw[ultra thick, brandPrimary] (0,0) -- (4,2) node[right] {2};
    
    % Точки
    \fill[black] (2,2) circle (1.5pt);
    \fill[black] (4,4) circle (1.5pt);
    \fill[black] (2,1) circle (1.5pt);
    \fill[black] (4,2) circle (1.5pt);
    
    % Пунктиры
    \draw[dashed, gray] (2,0) -- (2,2);
    \draw[dashed, gray] (4,0) -- (4,4);
    \draw[dashed, gray] (0,1) -- (2,1);
    \draw[dashed, gray] (0,2) -- (4,2);
    \draw[dashed, gray] (0,4) -- (4,4);
    
    \node[left] at (0,1) {1};
    \node[left] at (0,2) {2};
    \node[left] at (0,4) {4};
    \node[below] at (2,0) {2};
    \node[below] at (4,0) {4};
    \node[below left] at (0,0) {0};
\end{tikzpicture}
\end{center}

1) Сопротивление первого проводника больше сопротивления второго проводника.\\
2) Электрическое сопротивление второго проводника в $2$ раза больше сопротивления первого проводника.\\
3) Удельное сопротивление материала первого проводника больше удельного сопротивления материала второго проводника.\\
4) При одинаковом напряжении сила тока в первом проводнике в $2$ раза больше, чем во втором.\\
5) При одинаковой силе тока напряжение на концах первого проводника больше, чем на концах второго.

\kim{15}
\end{taskbasic}
\answer{24}

% === ЗАДАЧА 4 ===
% Тег: #ЭЛЕКТРОДИНАМИКА #ТОК #ЗАКОН_ОМА #КИМ_15 #ВАРИАНТ_18
\begin{taskbasic}[code={\label{task:elem_ohm_v18_n15}}]
\textbf{Задача 4.} На рисунке приведён график зависимости силы тока $I$ в реостате от напряжения $U$ на его зажимах. Из приведённого ниже списка выберите все верные утверждения относительно характеристик реостата.

\nopagebreak\smallskip
\begin{center}
\begin{tikzpicture}[scale=1.0, every node/.style={font=\footnotesize}]
    \draw[xstep=1, ystep=1, gray!30, thin] (0,0) grid (4,3);
    
    \draw[-{Stealth[scale=0.8]}, black, thick] (0,0) -- (4.5,0) node[right] {$U$, В};
    \draw[-{Stealth[scale=0.8]}, black, thick] (0,0) -- (0,3.5) node[above] {$I$, А};
    
    % График VAH (проходит через (0,0) и (3,2))
    \draw[ultra thick, brandPrimary] (0,0) -- (4.5,3);
    
    \fill[black] (3,2) circle (1.5pt);
    \draw[dashed, gray] (3,0) -- (3,2);
    \draw[dashed, gray] (0,2) -- (3,2);
    
    \node[left] at (0,2) {2};
    \node[below] at (3,0) {6};
    \node[below left] at (0,0) {0};
\end{tikzpicture}
\end{center}

1) Электрическое сопротивление реостата равно $3\text{ Ом}$.\\
2) При напряжении $3\text{ В}$ сила тока в реостате равна $2\text{ А}$.\\
3) С увеличением напряжения сопротивление реостата увеличивается.\\
4) При силе тока $1\text{ А}$ напряжение на реостате равно $1{,}5\text{ В}$.\\
5) При напряжении $12\text{ В}$ сила тока в реостате равна $4\text{ А}$.

\kim{15}
\end{taskbasic}
\answer{15}
% ====================================================================
% ФАЙЛ: tasks/03_electrodynamics/ohm_law_and_resistance_part3.tex
% ТЕМА: ЗАКОН ОМА ДЛЯ УЧАСТКА ЦЕПИ. УДЕЛЬНОЕ СОПРОТИВЛЕНИЕ
% ПАЧКА 3: ВАРИАНТЫ 21–30
% ====================================================================

% === ЗАДАЧА 1 ===
% Тег: #ЭЛЕКТРОДИНАМИКА #ТОК #ЗАКОН_ОМА #КИМ_11 #ВАРИАНТ_23
\begin{taskbasic}[code={\label{task:elem_ohm_v23_n11}}]
\textbf{Задача 1.} Во сколько раз уменьшится электрическое сопротивление цилиндрического металлического проводника, если его длину уменьшить в $4$ раза, а площадь поперечного сечения увеличить в $4$ раза?

\kim{11}
\end{taskbasic}
\answer{16}

% === ЗАДАЧА 2 ===
% Тег: #ЭЛЕКТРОДИНАМИКА #ТОК #ЗАКОН_ОМА #КИМ_11 #ВАРИАНТ_24
\begin{taskbasic}[code={\label{task:elem_ohm_v24_n11}}]
\textbf{Задача 2.} Во сколько раз увеличится сила тока в цилиндрическом металлическом проводнике при постоянном напряжении на его концах, если площадь его поперечного сечения увеличить в $4$ раза?

\kim{11}
\end{taskbasic}
\answer{4}

% === ЗАДАЧА 3 ===
% Тег: #ЭЛЕКТРОДИНАМИКА #ТОК #ЗАКОН_ОМА #КИМ_15 #ВАРИАНТ_25
\begin{taskbasic}[code={\label{task:elem_ohm_v25_n15}}]
\textbf{Задача 3.} На рисунке представлены графики зависимости силы тока $I$ от напряжения $U$ для двух металлических проводников 1 и 2, изготовленных из одинакового материала. Длина первого проводника в $2$ раза больше длины второго. Из приведённого ниже списка выберите все верные утверждения относительно этих проводников.

\nopagebreak\smallskip
\begin{center}
\begin{tikzpicture}[scale=1.0, every node/.style={font=\footnotesize}]
    \draw[xstep=1, ystep=1, gray!30, thin] (0,0) grid (4,4);
    
    \draw[-{Stealth[scale=0.8]}, black, thick] (0,0) -- (4.5,0) node[right] {$U$, В};
    \draw[-{Stealth[scale=0.8]}, black, thick] (0,0) -- (0,4.5) node[above] {$I$, А};
    
    % График 1: R1 = 2 Ом (I = U/2)
    \draw[ultra thick, brandPrimary] (0,0) -- (4,2) node[right] {1};
    
    % График 2: R2 = 0.5 Ом (I = 2U)
    \draw[ultra thick, brandPrimary] (0,0) -- (2,4) node[above] {2};
    
    % Точки
    \fill[black] (4,2) circle (1.5pt);
    \fill[black] (2,1) circle (1.5pt);
    \fill[black] (2,4) circle (1.5pt);
    \fill[black] (1,2) circle (1.5pt);
    
    % Пунктиры
    \draw[dashed, gray] (2,0) -- (2,4);
    \draw[dashed, gray] (4,0) -- (4,2);
    \draw[dashed, gray] (0,2) -- (4,2);
    \draw[dashed, gray] (0,4) -- (2,4);
    
    \node[left] at (0,2) {2};
    \node[left] at (0,4) {4};
    \node[below] at (2,0) {2};
    \node[below] at (4,0) {4};
    \node[below left] at (0,0) {0};
\end{tikzpicture}
\end{center}

1) Сопротивление первого проводника меньше сопротивления второго проводника.\\
2) Электрическое сопротивление первого проводника в $4$ раза больше сопротивления второго проводника.\\
3) Площадь поперечного сечения первого проводника меньше площади сечения второго в $2$ раза.\\
4) Удельное сопротивление материала первого проводника больше удельного сопротивления второго.\\
5) Площадь поперечного сечения первого проводника в $2$ раза больше площади сечения второго.

\kim{15}
\end{taskbasic}
\answer{25}

% === ЗАДАЧА 4 ===
% Тег: #ЭЛЕКТРОДИНАМИКА #ТОК #ЗАКОН_ОМА #КИМ_25 #ВАРИАНТ_26
\begin{taskbasic}[code={\label{task:elem_ohm_v26_n25}}]
\textbf{Задача 4.} Никелиновый проводник длиной $l = 10\text{ м}$ и площадью поперечного сечения $S = 0{,}4\text{ мм}^2$ подключён к источнику постоянного напряжения. Сила тока в проводнике равна $I = 1{,}2\text{ А}$. Чему равно напряжение на концах этого проводника? (Удельное сопротивление никелина $\rho = 0{,}4\text{ Ом}\cdot\text{мм}^2/\text{м}$). Ответ выразите в вольтах (В).

\kim{25}
\end{taskbasic}
\answer{12}